% Français 9115, batch 4 — Gallica views 61–80.
% Pass: Opus 5 (claude-opus-5), 2026-09-12 — first pass, unchecked against
% the views by a human.
%
% What the twenty views hold. One piece, running without a break: the
% « Mémoire sur la courbure des surfaces » that batch 3 opened at view 43 and
% left in mid-sentence at view 60. Every one of the twenty views carries
% text — no blank verso, no guard leaf, nothing skipped — and the batch again
% ends in mid-sentence, at the foot of view 80. The author's own pagination
% runs 17 to 36 across the twenty leaves, one page per view, which is the
% plainest proof that nothing is missing here.
%
% The argument of these twenty pages. Batch 3 ended inside N° 3, the law by
% which curvature is distributed around a point, argued alongside Fourier's
% law for the heated ring. Views 61 to 63 finish that comparison and then
% stand back from it; view 64 opens N° 4 and the memoir changes method. From
% here on she builds the thing the memoir is for: a construction that makes
% the quantity of curvature at a point visible as an area.
%
%   v. 61–62   End of N° 3. The angular distance between the two equivalent
%              positions of the normal plane, and its counterpart in the
%              ring — two points of equal temperature, a half circumference
%              apart where the curvatures are a quarter. She grants that
%              this last proposition is not in Fourier's text but follows
%              from his formulas. Then the sphere whose four quadrants lie
%              alternately inside and outside the surface.
%   v. 63–64   Why the analogy does not surprise her: the characteristic
%              notions of things are infinite in number, the necessary order
%              among their parts is not. Mean quantities in particular
%              « forme le lien entre l'uniformité et la variété », and a
%              phenomenon developing around a centre gives, by substituting
%              space for time, « une construction de la périodicité ».
%   v. 64–66   N° 4, back to first notions. Curvature judged by how far an
%              arc departs from the tangent at its end; that perpendicular
%              named the « ligne de distance », equal to the versed sine.
%              Two arcs of equal length on circles of radii $\rho m$ and
%              $\rho n$ subtend $\varphi n$ and $\varphi m$; passing to the
%              infinitely small, the lines of distance come out inversely
%              proportional to the radii. Hence the sentence the rest of the
%              memoir turns on: linear curvature may be said to be directly
%              proportional to the line of distance.
%   v. 66–70   N° 5, the same step taken for surfaces. Turn the normal plane
%              through every position, cut off arcs that would straighten
%              into equal lengths, join their ends: that curve is the
%              directrix, and the line of distance sweeping along it
%              generates the « surface des distances », to which the
%              curvature of the surface is proportional. Beside it the
%              « surface des distances moyennes », cylindrical, built on the
%              sphere of mean curvature. The two together separate what had
%              always been confounded — the figure of a surface and the
%              quantity of curvature it carries at a point.
%   v. 71–75   That the two have equal extent. First by the relations
%              already proved for the inverse radii, carried over to the
%              lines of distance; then, at length, by developing both
%              surfaces into the plane. The development is the parallelogram
%              $ABDC$ against the figure $ACDh'q'hqB$, and the equality
%              reduces to equal triangles, « tellement évidente qu'il seroit
%              inutile de nous arrêter à la démontrer ».
%   v. 75–77   N° 7, the propositions of N° 3 proved again from the
%              development, and not in their first order: 2nd prop. first
%              (v. 75), then 3rd (v. 75–76), then 1st (v. 76–77). With the
%              objection of N° 2 answered — she is not applying formulas of
%              linear curvature to surfaces, since a surface is compared to
%              its tangent plane exactly as a curve is to its tangent.
%   v. 78–80   The consequence she is after: $x$ running from zero to $mn$,
%              $pq = mn - x$ and $gh = mn + x$ give an infinity of figures,
%              from the sphere to the surface with one principal curvature
%              nil, all with the same quantity of curvature. Figure depends
%              on the difference of the principal curvatures, mean curvature
%              on their sum. N° 8 opens at view 80 and the batch stops
%              inside its first argument.
%
% The numbering on the leaves, and why no \folio{} is written. The ink
% series of batches 1 to 3 continues without a gap: 30, 31, 32, 33, 34, 35,
% 36, 37, 38, 39 on the rectos — views 61, 63, 65, 67, 69, 71, 73, 75, 77,
% 79. Beside each of those, in the same ink, stands the author's own
% pagination, odd and struck through: 17, 19, 21, 23, 25, 27, 29, 31, 33,
% 35. The even pagination stands unstruck in the top left corner of each
% verso — 18, 20, 22, 24, 26, 28, 30, 32, 34, 36 — and is not struck. Every
% one of these numbers was read on the view; none is inferred. They are in
% ink. No pencilled foliation of the Bibliothèque was found anywhere in the
% batch, so, as in batches 1 to 3, no \folio{} is written. The only pencil
% on these twenty views is a later reader's: an asterisk and a word in the
% margin of view 67, and a word between parentheses in the margin of view
% 78, neither read.
%
% Notation. $\rho$ is the element of radii and $\varphi$ the element of
% arcs, introduced at view 65; $m$ and $n$ are numbers with $m > n$, so that
% the radii are $\rho m$, $\rho n$ and the arcs $\varphi n$, $\varphi m$ —
% the crossing is deliberate, since the two arcs are to have equal length.
% $D$ and $d$ are the lines of distance of the more and the less curved arc.
% From view 73 the letters are those of a figure: $ABDC$, $mn$, $pq$, $gh$
% and their primes. She writes « fig. 1 » at view 73 and « fig 1 et 2 » at
% view 78, but no figure is drawn on any leaf of this batch. Two places
% where the leaf does not say what the sense wants are transcribed as she
% wrote them and noted at the spot: « l'arc le plus courbe $m\rho$ » at view
% 65, and the argument $(\varphi m)$ in the second pair of equations at view
% 66, where the development that follows requires $(\varphi n)$. Neither was
% corrected. « Il s'agit à démontrer » at view 74 is likewise hers. The
% parallelogram is $ABDC$ at views 73 and 80 and $ABCD$ at view 74.
%
% Her hand, and two things it does. Almost every word ends in a heavy
% terminal flourish that can be taken for a final « s »; readings here
% follow the sense where the stroke alone will not decide. Her tall
% unlooped « s » and her looped « l » are distinct, which is what settles
% « de signes opposés » at view 63. Her own underlining is set \emph{};
% \uncertain{} renders as an underline too, so a reader of the PDF should
% take this header's word for which is which. In the quotation from Euler
% at view 80 she repeats a guillemet at the head of every line; as in
% batches 2 and 3 only the opening and closing ones are kept here. The
% square brackets at views 69, 71, 72, 73 and 76 are hers, drawn on the
% leaf.
%
% What is corrected on the leaves. Views 69, 71, 72, 77 and 78 are worked
% over: words struck, replacements written between the lines, and at view 69
% a long addition that begins between two lines and then runs from top to
% bottom down the left margin. Each is transcribed where it stands and
% described in a \note{} at the spot. View 72 carries the same correction as
% view 69 — « la courbe rentrante … comme un cercle parfait » — added in
% three places in the margin.
%
% What could not be read. Nine places, none guessed: the struck cluster
% after « l'arc » at view 62; the struck word after « de l'arc » at view 68;
% the last two lines of the marginal addition at view 69, cramped at the
% foot of the margin; the word in superscript and the abbreviation before
% « N° 2 et 3 » in the interlinear addition at view 71; the circled note of
% three words after « N° 5 » at view 72; the word under the pencil asterisk
% at view 67; the word struck under « accompagneront » at view 78; and the
% pencil word in the margin of that same view. The word overwritten by
% « assujettis » at view 67 is likewise lost under the correction.
%
% Nothing here was checked against any published transcription.
\documentclass[11pt,a4paper]{article}
% The preamble shared by every transcription of the mathematicians' manuscripts in Gallica.
%
% It rests on one idea: **uncertainty compounds, it does not resolve.** A
% transcription that smooths over a doubtful word has destroyed the only thing
% it was made for — a reader must be able to tell, at every point, what was on
% the page and what was supplied. The apparatus macros below are therefore the
% critical apparatus, not decoration.
%
% The same commands are recognised by `scripts/render.mjs`, which produces the
% site's reading view, and by `scripts/tei.mjs`. A macro added here must be
% added there too, or rendering fails loudly — which is the wanted behaviour.
%
% Comments here are English, like the rest of the repository. The documents
% this preamble sets are French, and that is deliberate: see the skills.

\ProvidesPackage{germain}[2026/09/15 Transcriptions of mathematicians' manuscripts in Gallica]

\usepackage{amsmath,amssymb,amsthm}
\usepackage{mathrsfs}
% Kept from the parent preamble so the renderer's diagram support stays
% available; these manuscripts draw no commutative diagrams, and nothing here uses it.
\usepackage{tikz-cd}
\usepackage[english,french]{babel}
\usepackage{fontspec}

% --- Greek ------------------------------------------------------------------
% The text face stays fontspec's default, Latin Modern, which has no Greek:
% XeTeX drops every glyph it lacks with a line in the log and nothing on the
% page, and the Greek words the manuscripts quote vanished from the PDFs. So
% the Greek and Greek Extended blocks get a XeTeX character class of their own,
% and entering or leaving a run of it switches the family to CMU Serif — the
% same Computer Modern design, with polytonic Greek — and back. Only the
% family changes: a Greek word in \textit or \textbf keeps its shape and series.
% The transitions to and from every other class carry the tokens that class
% has towards an ordinary letter, so French spacing before « ; » or inside
% « » still holds after a Greek word. No grouping, so no brace or \endgroup
% can come out unbalanced; maths is left alone, since XeTeX inserts nothing
% there. Kept identical in `exercices.sty`.
\makeatletter
\newfontfamily\germainGreekfont{cmun}[
  Extension=.otf, NFSSFamily=germaingreek,
  UprightFont=*rm, ItalicFont=*ti, SlantedFont=*sl,
  BoldFont=*bx, BoldItalicFont=*bi, BoldSlantedFont=*bl]
\newXeTeXintercharclass\germain@greek
\def\germain@greekrange#1#2{%
  \@tempcnta=#1\relax
  \loop
    \XeTeXcharclass\@tempcnta=\germain@greek
    \ifnum\@tempcnta<#2\advance\@tempcnta\@ne\repeat}
\germain@greekrange{"0370}{"03FF}
\germain@greekrange{"1F00}{"1FFF}
\def\germain@greekprev{\rmdefault}
\def\germain@greekin{%
  \edef\germain@greekprev{\f@family}\fontfamily{germaingreek}\selectfont}
\def\germain@greekout{\fontfamily{\germain@greekprev}\selectfont}
\def\germain@greekjoin{%
  \XeTeXinterchartokenstate=\@ne
  \@tempcnta=\z@
  \loop
    \ifnum\@tempcnta=\germain@greek\else
      \XeTeXinterchartoks\germain@greek\@tempcnta=\expandafter{%
        \expandafter\germain@greekout\the\XeTeXinterchartoks\z@\@tempcnta}%
      \XeTeXinterchartoks\@tempcnta\germain@greek=\expandafter{%
        \the\XeTeXinterchartoks\@tempcnta\z@\germain@greekin}%
    \fi
    \ifnum\@tempcnta<4095\advance\@tempcnta\@ne\repeat}
% babel-french sets its own classes, and saves and restores the interchar
% state, each time a language is selected: join again after it does.
\addto\extrasfrench{\germain@greekjoin}
\addto\extrasenglish{\germain@greekjoin}
\AtBeginDocument{\germain@greekjoin}
\makeatother

\usepackage{geometry}
\geometry{a4paper,margin=25mm,marginparwidth=28mm}
\usepackage{marginnote}
\usepackage{xcolor}
\usepackage[normalem]{ulem}
\setlength{\emergencystretch}{3em}
\usepackage{hyperref}
\hypersetup{colorlinks=true,linkcolor=black,urlcolor=black,pdfborder={0 0 0}}

% --- Batch metadata ---------------------------------------------------------
% Taken as they stand by the rendering script, which shows them at the head of
% the reading view. They fix what the file claims to cover. The macro names are
% the parent project's, so that the scripts stay shared; \folder here means the
% volume's slug — `fr-9115` — and \pages counts Gallica views.

\newcommand{\folder}[1]{\gdef\grothFolder{#1}}
\newcommand{\batch}[1]{\gdef\grothBatch{#1}}
\newcommand{\pages}[2]{\gdef\grothFirst{#1}\gdef\grothLast{#2}}
\newcommand{\foldertitle}[1]{\gdef\grothFoldertitle{#1}}

% The holder's own shelfmark, printed as they write it: « Français 9115 ».
\newcommand{\shelfmark}[1]{\gdef\grothShelfmark{#1}}

% The dating, copied verbatim from the holder's catalogue — never inferred
% here. « XIXe siècle » is the BnF's; a pass that dates a leaf from its content
% says so in a \note{}, as its own inference.
\newcommand{\dating}[1]{\gdef\grothDating{#1}}

% The status notice, printed in the title block and repeated as a diagonal
% watermark on every page of the PDF. The manuscripts are in the public
% domain; what this line declares is not a right but a quality — an
% unverified first pass by a machine. The skills write the line; a file
% without it gets no watermark, so leaving it out is visible in the output.
\newcommand{\watermark}[1]{\gdef\grothWatermark{#1}}

\gdef\grothFolder{?}\gdef\grothBatch{}\gdef\grothFirst{?}\gdef\grothLast{?}%
\gdef\grothFoldertitle{}\gdef\grothShelfmark{}\gdef\grothDating{}\gdef\grothWatermark{}

% --- The résumé -------------------------------------------------------------
\newenvironment{resume}
  {\small\setlength{\parskip}{0.45em}}
  {\par\medskip\hrule\medskip}

% --- Keywords ---------------------------------------------------------------
% In English, closing the modernised reading's résumé: the modern vocabulary
% under which to search for what the leaves construct. The manifest extracts
% them to tag the volume — this line is the single source of the tags.
\newcommand{\keywords}[1]{\par\smallskip\noindent
  {\footnotesize\textsc{Keywords} --- \textit{#1}}\par}

% --- The critical apparatus -------------------------------------------------

% The Gallica view number, in the margin. This is the anchor that lets a
% disagreement over a formula be settled by looking at *one* image rather than
% a volume of 750. The site uses it to turn the facsimile as the transcription
% is scrolled. A view is one scanned image: usually one side of a leaf.
\newcommand{\page}[1]{%
  \marginnote{\footnotesize\color{gray!70}\textsf{v.\,#1}}%
  \label{page:#1}\ignorespaces}

% The BnF's pencilled foliation, where the leaf carries it and a pass can read
% it: « 198r ». This is what the literature cites — « ff. 198r–208v » — and
% the one line that turns a citation into a view. Recorded, never inferred:
% a folio a pass cannot read is not written.
\newcommand{\folio}[1]{%
  \marginnote{\footnotesize\color{gray!50}\textsf{f.\,#1}}\ignorespaces}

% The modernised reading's coarser counterpart to \page{}: which views a
% section is drawn from.
\newcommand{\pagerange}[2]{%
  \marginnote{\footnotesize\color{gray!70}\textsf{v.\,#1--#2}}%
  \ignorespaces}

% Illegible: never guessed.
\newcommand{\ill}{\textcolor{red!60!black}{[\,\ldots\,]}}

% A doubtful reading: the word is given, and flagged as uncertain.
\newcommand{\uncertain}[1]{\uline{#1}}

% An editorial addition — a missing letter, an implied word.
\newcommand{\add}[1]{\textcolor{blue!50!black}{[#1]}}

% Struck out by the writer. What was crossed out is part of the document.
\newcommand{\struck}[1]{\sout{#1}}

% The transcriber's note, on the state of the page or the mathematics.
\newcommand{\note}[1]{\par\smallskip{\small\color{gray!60!black}\textit{[#1]}}\par\smallskip}

% A marginal note of hers, distinct from the body of the text.
\newcommand{\marginal}[1]{\par\smallskip\noindent\hspace{1em}%
  {\small\color{gray!60!black}#1}\par\smallskip}

% --- Title ------------------------------------------------------------------
% The title block is French, like the documents themselves.

\AddToHook{shipout/background}{%
  \ifx\grothWatermark\empty\else
    \begin{tikzpicture}[remember picture, overlay]
      \node[rotate=52, scale=3.4, align=center, text=gray!18,
            font=\sffamily\bfseries]
        at (current page.center) {\grothWatermark};
    \end{tikzpicture}%
  \fi}

\AtBeginDocument{%
  \begin{center}
    {\large\bfseries Manuscrits de mathématiciens (Gallica) ---
      \ifx\grothShelfmark\empty\grothFolder\else\grothShelfmark\fi}\\[2pt]
    {\itshape\grothFoldertitle}\\[4pt]
    {\small vues \grothFirst--\grothLast
      \ifx\grothBatch\empty\else\ (lot \grothBatch)\fi}\\[2pt]
    \ifx\grothDating\empty\else
      {\small\color{gray!60!black}Datation du catalogue : \grothDating}\\[2pt]
    \fi
    \ifx\grothWatermark\empty\else
      {\small\bfseries\color{red!45!gray}\grothWatermark}\\[2pt]
    \fi
    {\footnotesize\color{gray!60!black}Transcription établie d'après les images IIIF de Gallica.
     Source gallica.bnf.fr / Bibliothèque nationale de France.\\
     Les lectures incertaines sont soulignées, les illisibles notées [\,\ldots\,],
     les ajouts entre crochets ; \textsf{v.} numérote les vues de Gallica,
     \textsf{f.} la foliotation au crayon de la BnF.}
  \end{center}
  \vspace{1em}}

\endinput


\folder{fr-9115}
\shelfmark{Français 9115}
\batch{4}
\pages{61}{80}
\dating{1801-1900}
\watermark{Lecture automatique\\non vérifiée}
\foldertitle{Papiers de Sophie GERMAIN. — Recueil de dissertations et problèmes mathématiques et physiques.}

\begin{document}

\note{Numérotation des feuillets. La série à l'encre des lots précédents se
poursuit sans interruption sur les rectos — vues 61, 63, 65, 67, 69, 71, 73,
75, 77, 79 → 30, 31, 32, 33, 34, 35, 36, 37, 38, 39. À côté de chacun de ces
nombres, de la même encre, la pagination de l'auteur, impaire et biffée :
17, 19, 21, 23, 25, 27, 29, 31, 33, 35. La pagination paire est portée, non
biffée, en haut à gauche de chaque verso : 18, 20, 22, 24, 26, 28, 30, 32,
34, 36. Tous ces nombres ont été lus sur la vue ; aucun n'est déduit. Ils
sont à l'encre : aucune foliotation au crayon de la Bibliothèque n'a été lue
sur ces vues, et aucun « f. » n'est donc porté dans cette transcription.}

\page{61}
contenues dans les différens plans normaux ; tandis que la position des
plans de courbures principales est unique, toute autre position du plan
normal est équivalente à une seconde position du même plan.

L'angle compris entre les deux positions équivalentes du plan normal est
évidemment égal à $\frac{\pi}{2} - 2\varphi'$, différence entre les deux
valeurs assignées à l'angle $\varphi$.

Un arrangement analogue se manifeste par rapport à la distribution de la
chaleur dans l'anneau ; car toute température comprise entre les
températures extrèmes appartient à la fois à deux points différens de
l'anneau, et celle-ci ne convient qu'à un seul.

Cette proposition ne peut, à la vérité, être, comme les précédentes,
extraite textuellement du livre cité ; mais elle résulte des formules qui y
sont données.

La position des deux points dont la température est semblable sera
représentée comme il suit : $\varphi$ étant l'arc compris entre un de ces
points, choisi arbitrairement, et celui auquel appartient la température la
plus élevée ; $\varphi'$ l'arc entre le même point choisi et le point dont
la température est en même tems la température moyenne de l'anneau, on aura
$\varphi = \frac{\pi}{2} + \varphi'$ ; et, à l'égard du point dont la
température est égale à celle du point choisi,
$\varphi = \frac{3\pi}{2} - \varphi'$.

L'arc compris entre les deux points dont la température est semblable sera
donc $\pi - 2\varphi'$.

Si on a présente à la pensée la différence que nous avons

\page{62}
signalée relativement aux distances angulaires qui, dans les lois que nous
comparons, séparent les quantités extrèmes ; on verra une analogie sensible
résulter des rapprochemens suivans.

D'une part, il existe, en général, quatre, et, pour un cas particulier, deux
courbes semblables : de l'autre, en général deux, et, pour un cas
particulier, un seul point de l'anneau, jouissent d'une même température :

La distance angulaire qui sépare soit les courbes semblables, soit les
points d'égale température, est exprimée par la différence entre la distance
angulaire qui sépare les valeurs extrèmes, et l'arc \ill{} compris entre le
lieu des valeurs moyennes et la courbe ou le point choisi.

Une conséquence immédiate de cette dernière proposition nous ramène à ce
qu'on sait déjà touchant les sphères qui ont leur centre sur la normale à la
surface.

Il est évident, en effet, que, les courbures principales n'étant contenues
que dans un seul des plans normaux, les sphères décrites de chacun de leur
rayon ne couperont pas la surface ; tandis que, par une conséquence inverse,
les sphères décrites des rayons compris entre ces deux limites couperont au
contraire la même surface.

Parmi les sphères qui coupent la surface, il convient de distinguer celle
dont les quatre quadrans sont situés alternativement au dedans et au dehors
de la surface. Le rayon de cette sphère appartient en même tems à la courbe
contenue dans le plan normal qui fait l'angle de $45^{\circ}$ avec ceux des
courbures principales, et la proposition 2\textsuperscript{ème} nous apprend que
la courbure contenue dans le plan

\page{63}
dont il s'agit est moyenne entre celles-ci. La courbure des grands cercles
de notre sphère est donc moyenne entre les courbures principales de la
surface.

Nous savons aussi que les deux plans normaux qui contiennent les courbures
moyennes partagent le contour du point donné en quatre parties, dont l'état
est pareil, mais alternativement de signes opposés ; et il est clair
d'ailleurs que ces plans coupent la sphère en même tems que la surface.
L'intersection de la sphère et de la surface partage donc aussi le contour
du point donné en quatre parties, dont l'état est pareil, mais de signes
opposés ; et la situation de ces parties, les unes intérieures, les autres
extérieures à la sphère, rend sensible cette répartition de la courbure
autour du point donné.

L'analogie que nous avons fait remarquer entre deux lois applicables à des
sujets entièrement dissemblables n'étonnera point, ne sait-on pas, en effet,
que, si les notions caractéristiques de chacun des objets de nos
connoissances sont en nombre infini, l'ordre nécessaire entre les diverses
parties de ces objets ne sauroit présenter la même variété ? Habitués à
reconnoître, de quelque côté que nous jetions les yeux, qu'un petit nombre
de lois, parmi lesquelles les plus simples sont aussi celles dont
l'application a lieu le plus souvent, se partagent l'immense domaine de la
réalité, il nous seroit également impossible de concevoir un ordre d'idées
ou de faits qui puisse se

\page{64}
soustraire à la nécessité d'un ordre déterminé et de comprendre comment un
ordre marqué au sceau de la simplicité, s'associeroit au développement d'une
idée erronnée.

L'existence des quantités moyennes, surtout, se fait sentir dans une
multitude de considérations de natures différentes ; elle forme le lien
entre l'uniformité et la variété. Si le phénomène qui en présente l'exemple
se développe autour d'un point pris pour centre, la représentation des
valeurs diverses observées à la fois offrira, en quelque sorte, par la
substitution de l'espace au tems, une construction de la périodicité.

Après avoir cherché à étayer la considération des moyennes courbures sur une
analogie que des rapports faciles à saisir auroient pû faire prévoir, nous
ne craindrons plus de traiter isolément une question dont l'existence, ainsi
établie d'avance, ne doit laisser aucun doute dans l'esprit du lecteur.

N° 4 Commençons, ici, avant d'entrer en matière par remonter jusqu'aux notions
les plus communes touchant la courbure.

D'abord à l'égard des courbes linéaires.

L'idée qui se présente la première est de comparer la courbe à sa tangente ;
ainsi, après avoir fait choix de deux lignes droites d'égale longueur, si on
en forme les arcs de deux cercles différens, et qu'on mène leurs tangentes à
un des points d'extrémité de ces arcs, on jugera que celui des deux dont la
seconde extrémité

\page{65}
s'éloigne le plus de la tangente est aussi celui qui a le plus de courbure.

La distance entre l'extrémité d'un arc et la tangente menée à l'autre
extrémité du même arc est mesurée par une perpendiculaire à la tangente.

Nous nommerons dorénavant cette perpendiculaire \emph{ligne de distance}.

La considération des \emph{lignes de distances} devant être, dans la suite,
d'une grande utilité, nous nous arrêterons un instant à examiner leur
expression.

Il est évident que la \emph{ligne de distance} est égale au sinus verse de
l'angle correspondant ; car cette ligne est parallèle au rayon mené au point
de tangence, et elle est comprise entre la tangente et le sinus.

Cela posé, $\rho$ étant l'élément des rayons, $\varphi$ celui des arcs, $m$
et $n$ des nombres tels que les rayons des cercles auxquels appartiennent
les arcs dont nous comparons la courbure soient $\rho m$ et $\rho n$, ces
arcs eux mêmes seront $\varphi n$ et $\varphi m$.

On sait que la courbure du cercle est en raison inverse de son rayon ; en
prenant donc, pour fixer les idées, $m > n$, il est certain que le plus
grand rayon $m\rho$ et l'arc le moins courbe $n\varphi$ appartiendront à un
même cercle. Par la même raison, le plus petit rayon $n\rho$ et l'arc le
plus courbe $m\rho$ appartiennent aussi au même cercle.

\note{la vue porte bien $m\rho$ pour le dernier de ces quatre symboles, là
où la définition qui précède — les arcs sont $\varphi n$ et $\varphi m$ —
demande $m\varphi$. Le $\rho$ et le $\varphi$ sont ici deux tracés
nettement distincts, et rien n'a été corrigé. Une petite croix est portée
sous la ligne après $n\varphi$, sans renvoi lisible ailleurs sur le
feuillet.}

Soient $D$ et $d$ les \emph{lignes de distances} relatives aux arcs le plus
courbe et le moins courbe,

\page{66}
on aura, conformément à ce qu'on vient de dire,
\begin{align*}
  D &= n\rho \sin.\text{vers.}(\varphi m)\\
    &= n\rho\,(1 - \cos(\varphi m))\\
    &= n\rho \cdot 2\sin^{2}\tfrac{1}{2}(\varphi m) ,
\end{align*}
\begin{align*}
  d &= m\rho\,(1 - \cos(\varphi m))\\
    &= m\rho \cdot 2\sin^{2}\tfrac{1}{2}(\varphi m) ;
\end{align*}

\note{les deux colonnes sont écrites côte à côte sur le feuillet. L'argument
des deux dernières lignes est bien $(\varphi m)$ sur la vue, là où la
définition de la page précédente — le rayon $m\rho$ va avec l'arc
$n\varphi$ — demande $(\varphi n)$ ; le développement qui suit,
$d = \frac{1}{2}m.n^{2}\rho.\varphi^{2}$, le confirme. Rien n'a été corrigé
ici.}

si une courbe se confond avec son cercle osculateur, ce n'est que dans
l'étendue d'un arc infiniment petit, et nous ne nous occupons ici des
\emph{lignes de distances} que pour en faire l'application aux courbes.
Prenant donc $\varphi m$ et $\varphi n$ infiniment petits, nous trouverons
\begin{align*}
  D &= \tfrac{1}{2}\,n.m^{2}\rho.\varphi^{2}\\
    &= \frac{1}{n\rho}\cdot\tfrac{1}{2}m^{2}n^{2}\rho^{2}\varphi^{2} ,
\end{align*}
\begin{align*}
  d &= \tfrac{1}{2}\,m.n^{2}\rho.\varphi^{2}\\
    &= \frac{1}{m\rho}\cdot\tfrac{1}{2}m^{2}n^{2}\rho^{2}\varphi^{2} :
\end{align*}

Les \emph{lignes de distances} sont donc en raison inverse des rayons de
courbures, et, par conséquent, on peut dire également de la courbure
linéaire, ou, qu'elle est en raison inverse des rayons de courbures, ou,
qu'elle est en raison directe des \emph{lignes de distances}.

Lorsqu'on se borne à la considération des courbes, cette dernière manière de
s'exprimer seroit, sans doute, à la fois bizarre et surabondante : mais nous
verrons tout à l'heure qu'elle présente de grands avantages quand on cherche
à se former une idée complète de la courbure des surfaces.

N° 5 Si nous jugeons qu'une ligne est courbe lorsque nous la voyons s'écarter de
la direction d'une droite qui la touche, l'idée de la courbure des surfaces
nous est également suggérée par l'observation de la distance

\page{67}
qui sépare le plan tangent des points de la surface voisins du plan de
tangence.

Pour toute autre surface que celle de la sphère, cette distance n'est pas
uniforme.

La question des courbures moyennes peut être énoncée en ces termes :
\emph{trouver l'expression de la distance moyenne entre le plan tangent et
les points de la surface qui environnent le point de tangence.}

\note{un astérisque au crayon est porté dans la marge de gauche à la hauteur
de cette phrase, et un second à la fin de la phrase même ; sous celui de la
marge, un mot au crayon, \ill{}. Le crayon n'est pas de la main de
l'auteur.}

Occupons nous d'abord des seules surfaces dont les rayons de principales
courbures sont dirigés du même côté.

On a vû plus haut que la courbure d'une simple courbe est proportionnelle à
une ligne droite, la \emph{ligne de distance} ; nous allons montrer que la
courbure des surfaces est proportionnelle à une surface engendrée par une
\emph{ligne de distance} qui change à chaque instant de grandeur et de
position.

Il s'agit présentement de décrire la courbe qui doit diriger le mouvement de
la génératrice.

Après avoir mené par le point choisi le plan normal à la surface, on donnera
successivement à ce plan toutes les positions dont il est susceptible. Sur
chacune des courbes produites par les diverses intersections du plan normal
et de la surface on prendra des arcs infiniment petits, assujettis à cette
condition que, s'ils étoient étendus en lignes droites, ils donneroient tous
des longueurs égales. La courbe qui joint les extrémités de ces arcs sera la
directrice que nous avons cherchée.

\note{« assujettis » est écrit par-dessus un premier mot que la surcharge
rend illisible.}

\page{68}
Cette courbe est à double courbure.

Si on conçoit le plan tangent à la surface dans le point donné et la
perpendiculaire à ce plan qui tombe sur un des points de la directrice il
est clair que cette ligne ne sera autre que la \emph{ligne de distance}
contenue dans celui des plans normaux qui coupe la directrice dans le point
choisi.

La \emph{ligne de distance}, en changeant à chaque instant de grandeur, et
en parcourant successivement tous les points de la courbe qui dirige son
mouvement, sans cesser d'être perpendiculaire au plan tangent, engendrera
une surface que les plans de moyennes courbures, c'est-à-dire ceux qui font
l'angle de $45^{\circ}$ avec les plans de principales courbures, diviseront
en quatre parties.

Il est clair que la surface dont nous nous occupons est composée des
\emph{lignes de distances}, en nombre infini, qui appartiennent aux courbes,
également en nombre infini, formées par l'intersection des différens plans
normaux qui passent par le point donné de la surface, point par rapport
auquel on veut connoître l'expression de la courbure. La surface composée
des \emph{lignes de distances} sera nommée, par cette raison, \emph{surface
des distances}.

On a vu, N° 4, que la courbure d'une simple ligne courbe est proportionnelle
à la \emph{ligne de distance} de l'arc \ill{} cercle osculateur, en nous
laissant conduire ici par l'analogie, nous dirons que, dans le point donné,

\note{après « de l'arc », un mot d'environ six lettres est biffé et
recouvert : il n'a pas été lu.}

\page{69}
la courbure de la surface est proportionnelle à la \emph{surface des
distances} décrite autour de ce point.

La \emph{surface des distances} est le segment d'une surface cylindrique
ayant pour base \struck{l'ellipse} la courbe rentrante que les positions
successives de la génératrice tracent sur le plan tangent ; ce segment est
terminé à l'autre extrémité par la directrice : \struck{sa forme} cette
courbe dépend de la figure de la surface à laquelle appartient la
\emph{surface des distances}.

\marginal{avant d'aller plus loin il est bon de faire remarquer que la
petite courbe rentrante qui sert de base à la \emph{surface des distances}
peut toujours être regardée comme un cercle parfait. En effet il est évident
que la plus grande différence qu'on puisse imaginer entre les deux diamètres
de cette courbe correspond au cas des surfaces cylindriques ; mais alors en
adoptant les suppositions et les dénominations de N° 1, l'un de ces
diamètres devient \emph{a} \struck{l'autre} \ill{}}

\note{cette addition est de la main de l'auteur : elle commence entre les
lignes, au-dessus de « remarquer que la petite courbe rentrante », et se
poursuit de haut en bas dans la marge de gauche. Ses deux dernières lignes,
serrées au bas de la marge, n'ont pas été lues.}

Les plans de courbures moyennes contiennent évidemment les \emph{lignes de
distances moyennes} ; ils coupent à la fois la \emph{surface des distances},
qu'ils partagent en quatre parties symétriques ; la surface donnée, et la
sphère de moyenne courbure, défini N° 2 [Dont les quatre quadrans sont
alternativement situés au dedans et au dehors de cette dernière surface.

On a vû N° 3, 3\textsuperscript{ème} prop., que la courbure des grands cercles de
la sphère dont nous parlons, est moyenne entre les courbures principales de
la surface donnée.

Nous allons démontrer présentement que la courbure de la sphère elle même
est égale à la courbure moyenne de la surface donnée. Cette sphère sera
nommée dorénavant \emph{sphère de moyenne courbure de la surface}]

\note{« de moyenne courbure, défini N° 2 » est ajouté entre les lignes. Le
crochet ouvrant est porté devant « Dont », le crochet fermant après
« surface » : l'auteur a encadré tout ce passage.}

N° 6 Lorsque la surface est sphérique, la courbure autour de chacun de ses points
est uniforme ; et, par conséquent, \emph{dans un rayon donné}, les points
environnans le point de tangence sont tous également distans de ce point.

\page{70}
décrit, sur la sphère de moyenne courbure, un cercle dont le rayon soit le
sinus de l'arc, ou, à cause de sa petitesse, l'arc même, à l'extrémité
duquel appartient la \emph{ligne de distance} ; ce cercle sera la directrice
du mouvement de la même \emph{ligne de distance}, qui, sans changer de
grandeur et sans cesser d'être perpendiculaire au plan tangent, engendrera
la surface à laquelle la courbure de la sphère est proportionnelle. Cette
surface sera celle des \emph{distances moyennes}.

\note{« à l'extrémité » est ajouté entre les lignes, au-dessus de
« duquel ».}

La \emph{surface des distances moyennes} est cylindrique ; elle a pour base
le cercle que les positions successives de la génératrice tracent sur le
plan tangent : elle est terminée, à l'autre extrémité, par la directrice.

Si, pour deux surfaces données, de figures différentes, la somme des
courbures principales, et, par conséquent aussi, le rayon de la sphère de
moyenne courbure, sont les mêmes, la \emph{surface des distances moyennes}
sera aussi la même.

Ainsi, tandis que la \emph{surface des distances} dépend à la fois et de la
courbure moyenne et de la figure de la surface donnée, la considération de
la \emph{surface des distances moyennes}, isolant l'idée de la figure de
celle de la distance au plan tangent, offre une représentation séparée
\struck{de la} des deux affections qui ont toujours été confondues ;
savoir, la figure et la quantité de courbure qu'une surface possède dans
chacun de ses points.

La justesse de ce raisonnement exige que les deux

\page{71}
\emph{surfaces des distances}, quoique différentes par leur forme, ayent
pourtant la même étendue, puisque de cette étendue dépend évidemment, la
quantité de courbure, ou, en d'autres termes, la somme des distances entre
le plan tangent et les points de la surface qui environnent le point de
tangence.

Pour établir l'égalité d'étendue entre les deux \emph{surfaces des
distances}, il suffiroit peut être de faire remarquer qu'il existe entre
\emph{les lignes de distances} les mêmes relations qu'entre les raisons
inverses des rayons de courbures.

\note{au-dessus de « les raisons inverses des rayons de courbures »,
l'auteur a ajouté entre les lignes : « courbures \struck{que} auxquelles ces
lignes sont proportionnelles ».}

[Ainsi on a vu, N° 2, que la somme des raisons inverses des rayons de
courbures contenues dans deux plans normaux perpendiculaires entr'eux est
une quantité indépendante du choix de ces plans ; et, plus loin, que la
raison inverse du rayon contenu dans le plan de courbure moyenne, est elle
même moyenne entre les deux premières.]

On dira donc, aussi, que la somme des \emph{lignes de distances} contenues
dans deux plans normaux perpendiculaires entr'eux est indépendante du choix
de ces plans, et que la \emph{ligne de distance} moyenne, c'est-à-dire la
génératrice de la \emph{surface des distances moyennes} est elle même
moyenne entre les deux premières \emph{lignes des distances}.

\note{une addition de l'auteur court sur deux interlignes, de « est elle
même moyenne entre les » à « On dira donc » : « conformément \ill{} la loi
de distribution de la courbure exposée \ill{} N° 2 et 3 ». Le mot porté en
exposant au-dessus de « de la loi » n'a pas été lu, non plus que
l'abréviation de deux ou trois signes qui précède « N° 2 et 3 ». Le crochet
ouvrant est devant « Ainsi », le fermant après « premières ».}

Cela posé, quand on prend, sur l'une et l'autre des \emph{surfaces des
distances}, un nombre de positions, convenablement

\page{72}
choisies, des génératrices, la somme de ces lignes est égale de part et
d'autres. Si le nombre des positions de la génératrice est infini, ces
positions épuiseront l'étendue de la surface. On en peut conclure que
l'étendue elle même est égale de part et d'autre.

L'égalité d'étendue entre les deux \emph{surfaces des distances} étant, ici,
d'une grande importance, nous allons essayer de la rendre sensible ; et,
pour cela, nous aurons recours à la considération du développement de ces
surfaces.

On a vu, N° 5, que les plans de courbures moyennes coupent à la fois : la
surface donnée ; la \emph{surface des distances}, qu'ils partagent en quatre
parties symétriques ; et la sphère de courbure moyenne. Les mêmes plans
coupent donc aussi la \emph{surface des distances moyennes}, qu'ils
partagent en quatre parties égales ; et, par conséquent il y a
intersection, en même tems, entre les deux \emph{surfaces des distances}.

\marginal{\ill{}}

\note{après « N° 5 », une note de l'auteur, entourée d'un trait : trois mots
dont le dernier est un numéro précédé de « N° ». Elle n'a pas été lue.}

Nous avons dit plus haut que la courbe rentrante qui sert de base à la
\emph{surface des distances} peut être regardée comme un cercle parfait.

\note{« rentrante », « comme un » et « cercle parfait » sont portés dans la
marge de gauche, en trois endroits, en regard des lignes où ils s'insèrent.
C'est la même correction que celle de la vue 69.}

Cela posé, soit ABDC le développement de la \emph{surface des distances
moyennes} ; nous supposerons que le développement de l'ellipse qui sert de
base au segment cylindrique dont se compose la \emph{surface des distances}
se confond avec AC, développement du cercle qui sert de base à \struck{la}
l'une et à l'autre des \emph{surfaces} \struck{des distances}.
[\emph{surface}, cylindrique, des \emph{distances moyennes} ; supposition
d'autant plus admissible que, ces deux périphéries, de même ordre que leurs
rayons, étant elles mêmes

\note{ce passage est retravaillé. « l'une et à » est ajouté au-dessus de
« à la » ; « des distances » est biffé et noyé d'encre ; un « sera le » est
porté entre les lignes au-dessus de « dont se compose » ; le crochet
ouvrant est devant « surface, cylindrique ». L'ordre dans lequel l'auteur
voulait que ces corrections se lisent n'est pas établi : les lignes sont
données ici dans l'ordre où elles courent sur le feuillet.}

\page{73}
infiniment petites, leur différence est tout à fait négligeable]

Le parallélogramme $ABDC$ sera divisé en quatre parties égales au moyen des
trois lignes, $mn$, $m'n'$ et $m''n''$ parallèles aux côtés $AB$ et $CD$ :
on divisera encore en deux parties égales les intervalles $Am$, $mm'$,
$m'm''$ et $m''C$ ; des points de division on portera, toujours
parallèlement aux côtés $AB$ et $CD$, les lignes $pq$, $p'q'$ de moindres,
et $gh$, $g'h'$ de plus grandes distances, disposées alternativement :
enfin on joindra par des droites les points $B$ et $q$ ; $q$ et $h$ ; $h$ et
$q'$.

\note{« fig. 1 » est porté entre les lignes au-dessus de « sera divisé ».
Aucune figure n'est dessinée sur les feuillets de ce lot.}

Nous établirons d'abord que la figure $ACDh'q'hqB$ n'est autre que le
développement de la \emph{surface des distances}.

On sait déjà que $AC$ est le développement de la base de cette surface, et
que les plans de distances moyennes coupent à la fois les deux
\emph{surfaces des distances} qu'ils partagent en quatre parties symétriques
par rapport à l'une, égales par rapport à l'autre ; par construction, les
lignes $mn$, $m'n'$, $m''n''$ et $AB = CD$ partagent le développement
$ABDC$ de la \emph{surface des distances moyennes} en quatre parties égales,
ces lignes sont donc le développement des quatre intersections des deux
\emph{surfaces des distances}.

Si on conçoit que l'une des lignes d'intersection devienne mobile, $AB$, par
exemple, parcoure tous les points de $AC$ sans cesser de lui être
perpendiculaire, et acquière en même tems les différentes valeurs comprises
entre celles des \emph{lignes de moyennes et de moindres distances} $AB$ et
$pq$, puis, entre les lignes \emph{de moindres et de plus grandes
distances} $pq$ et $gh$ \&c, la seconde des extrémités de la ligne mobile
tracera sur le

\page{74}
plan la ligne brisée $Bqhq'h'D$.

Il s'agit à démontrer que les valeurs successives de la ligne mobile qui
trace sur le plan la figure $ACDh'q'hqB$ et celles de la génératrice de la
\emph{surface des distances} sont égales, et se succèdent dans le même
ordre.

\note{« à démontrer » est bien ce que porte la vue, là où l'usage demande
« de démontrer ». Rien n'a été corrigé.}

Ces valeurs sont égales, puisque, de part et d'autre, elles reçoivent
successivement toutes celles qui sont comprises entre les mêmes limites,
savoir, les lignes de moindre et de plus grande distance. Elles se
succèdent de part et d'autre dans le même ordre, car les parties aliquotes
du développement, $AC$, de la base de la \emph{surface des distances},
correspondent évidemment aux mêmes parties aliquotes de cette base, ensorte
que, par exemple, les deux positions de la ligne mobile que sépareroit un
intervalle égal au quart de $AC$, seroient les développemens des
\emph{lignes de distances} contenues dans deux plans normaux
perpendiculaires entr'eux.

La ligne brisée $Bqhq'h'D$ qui sert de limite aux diverses grandeurs de la
ligne mobile parallèle à $AB$ est donc le développement de la directrice de
la \emph{surface des distances} ; et, par conséquent, la figure
$ACDh'q'hqB$ est en effet celui de cette surface elle même.

Il reste à prouver qu'il y a égalité d'étendue entre le parallélogramme
$ABCD$ et la figure $ACDh'q'hqB$ ; c'est-à-dire entre les développemens des
deux surfaces

\note{le parallélogramme est ici nommé $ABCD$, et $ABDC$ à la vue 73 et plus
bas ; les deux graphies sont celles des vues.}

\page{75}
des \emph{distances}. Cette égalité résulte de celle entre les triangles
semblables, $Bqn$, $nhn'$, \&c ; et celle-ci est tellement évidente qu'il
seroit inutile de nous arrêter à la démontrer.

N° 7 La considération du développement des \emph{surfaces}, \emph{des distances}
et des \emph{distances moyennes}, nous ouvre une voie nouvelle pour arriver
à la démonstration des propositions qui ont été traitées N° 3. Sans nous
astreindre à l'ordre dans lequel elles ont été présentées d'abord, nous
reproduirons ici ces diverses propositions.

2\textsuperscript{ème} prop. \emph{La courbure contenue dans le plan normal qui
fait l'angle de $45^{\circ}$ avec ceux de plus grande et de moindre courbure
est moyenne entre les courbures principales de la surface.}

Conformément à ce qu'on a vû plus haut, la première de ces courbure est
proportionnelle à la \emph{ligne de moyenne distance} $mn$, fig 1, et les
courbures principales sont également proportionnelles aux \emph{lignes}
$gh$, $pq$, \emph{de plus grande et de moindre distance}. Pour justifier la
présente proposition, il suffit donc d'établir que $mn$ est moyenne entre
$gh$ et $pq$.

Menons $qt$ égale et parallèle à $pq$ ; on a, par construction $pm = mq$ ;
d'où résulte $qs = st$ ; les triangles semblables $qsn$, $qtb$ ; donneront
donc $ht = 2sn$ ; et on aura, par conséquent,
\[ pq = mn - sn, \quad gh = mn + sn, \quad pq + gh = 2mn . \]

\note{les points $s$, $t$ et $b$ ne sont introduits nulle part ailleurs dans
le lot, et aucune figure ne l'accompagne ; « $qtb$ » et « $ht$ » sont bien
ce que porte la vue, le $b$ et le $h$ y étant deux tracés distincts.}

3\textsuperscript{ème} prop. \emph{Il existe, en général, autour du point donné
quatre directions suivant lesquelles les courbes contenues}

\page{76}
\emph{dans chacun des plans normaux sont semblables. À l'égard des courbures
principales, le nombre des courbes semblables se réduit à deux.}

En d'autres termes, \emph{il existe, en général, quatre lignes de distances
égales entr'elles ; mais les distances principales sont au nombre de deux.}

Il est clair, en effet, que les deux parallèles à $pq$, terminées d'une part
à $AC$ de l'autre à la ligne brisée $Bqn$, et les deux parallèles à $p'q'$
comprises entre les mêmes limites, si elles sont menées à égale distance de
$pq$ pour les uns, et de $p'q'$ pour les autres, seront égales ; il y a donc
quatre pareilles lignes ; mais, si la distance entre $pq$ et ses parallèles,
et de même entre $p'q'$ et ses parallèles, est nulle ; elles se confondront
avec $pq$ pour les uns, avec $p'q'$ pour les autres ; les \emph{lignes de
distances principales} seront donc au nombre de deux seulement.

1\textsuperscript{re} prop. \emph{Quelle que soit d'ailleurs la position de deux
plans normaux perpendiculaires entr'eux, la somme des courbures contenues
dans ces deux plans est une quantité constante ; cette somme est, par
conséquent, égale à celle des courbures principales de la surface.}

Ou, en adaptant les expressions dont l'équivalence a été démontrée.

\emph{Quelle que soit d'ailleurs la position de deux lignes de distances,
si, dans le développement de la surface des distances, l'intervalle qui
sépare ces lignes est égal à}

\page{77}
\emph{celui qui sépare les lignes de principales distances, leur somme sera
la même ; et, par conséquent, cette somme sera égale à celle des lignes de
principales distances ;} proposition de toute évidence ; car si on conçoit
que le système des \emph{lignes de principales distances} se meut en
parcourant tous les points de $AC$ et, par conséquent, aussi tous ceux de
$Bqnh$ \&c, il est visible que dans une quelconque des positions de ce
système, l'accroissement obtenu par une des lignes qui le composent sera
égal à la diminution que l'autre aura éprouvée.

La manière dont nous envisageons présentement la loi de la distribution de
la courbure autour d'un des points de la surface écarte entièrement
l'objection tirée de l'emploi exclusif, dans la question de la courbure des
surfaces, de formules propres à la courbure linéaire, objection signalée
N° 2.

\note{ce paragraphe est marqué d'un long trait courbe dans la marge de
gauche.}

Ici, au contraire, la manière dont on obtient la mesure des deux genres de
courbures est entièrement fondée sur la nature différente de chacune d'elles.
En effet, si, quand on dit de la courbure linéaire, qu'elle est
proportionnelle à la \emph{ligne de distance} on compare tacitement la
courbe à sa tangente ; lorsqu'on dit de la courbure d'une surface qu'elle
est proportionnelle à la \emph{surface des distances}, on compare également
la surface à son plan tangent.

On voit encore par là comment il arrive que l'idée de la quantité de
courbure qui, dans le cas linéaire, est inséparable de celle de la figure ;
\struck{doivent} en être distinguée à l'égard des surfaces ; car, la
\emph{ligne de distance} étant déterminée de grandeur et de position, elle
doit rester la même tant que

\note{« doivent » est biffé sur sa fin ; « et distincte » est ajouté entre
les lignes au-dessus de « en être ».}

\page{78}
la courbure ne change pas, et elle ne peut convenir, par conséquent qu'à une
seule courbe. mais l'étendue de la \emph{surface des distances} peut être
déterminée sans que la forme de cette surface soit fixée ; ensorte que des
figures différentes de la surface donnée ; accompagneront pourtant la même
\emph{surface des distances moyennes}.

\note{« accompagneront » est écrit entre les lignes au-dessus d'un mot biffé
et recouvert d'encre, qui n'a pas été lu. Un mot au crayon, entre
parenthèses, est porté dans la marge de gauche à la hauteur de cette
phrase : il n'a pas été lu non plus, et le crayon n'est pas de la main de
l'auteur.}

Le double développement de la \emph{surface des distances moyennes} et de
celle \emph{des distances}, représenté fig 1 et 2, va nous servir à suivre
des yeux le nombre infini des changemens de forme que peut subir la
\emph{surface des distances}, et, par conséquent aussi la figure de la
surface donnée ; sans que la surface \emph{des distances moyennes}, c'est à
dire la quantité de courbure en éprouve aucun.

Lorsque la surface donnée est sphérique, on a $pq = gh = mn$ ; par
conséquent les angles en $q$ et $q'$ disparaissent, et la \emph{surface des
distances} se confond avec celle \emph{des distances moyennes} ; ce qui doit
être, en effet, puisque la surface donnée elle même se confond avec la
sphère de moyenne courbure.

$x$ étant une grandeur linéaire comprise entre les limites zéro et $mn$, on
peut prendre $pq = mn - x$, $gh = mn + x$. à chacune des valeurs
particulières de $x$ correspondra une figure déterminée de la surface
donnée. Les angles en $h$ et $h'$ s'éloigneront d'autant plus de $BD$,
développement de la directrice de la \emph{surface des distances moyennes},

\page{79}
et les angles en $q$ et $q'$ s'approcheront d'autant plus de $AC$,
développement commun des bases des deux \emph{surfaces des distances}, que
la valeur de $x$ sera plus grande. Quand la grandeur de $x$ aura atteint sa
limite $mn$, les angles en $q$ et $q'$ auront leur sommet sur $AC$ ; ce qui
nous avertit qu'il y aura, alors, deux directions du plan normal pour
lesquelles la \emph{ligne de distance} sera nulle. Nous savons que, si on
trouve deux, et non pas quatre \emph{lignes de distances} égales entr'elles,
elles appartiennent à une des courbures principales de la surface. la
surface donnée à laquelle conviendront les développemens représentés fig. 2
aura donc une de ses courbures principales nulle.

On voit \struck{donc} qu'entre les figures sphériques et cylindrique de la
surface donnée, il existe un nombre infini de figures qui peuvent être
attribuées à cette surface, sans qu'il en résulte le moindre changement
relatif à la quantité de la courbure, quantité représentée ici par le
développement de la \emph{surface des distances moyennes}. C'est en effet ce
qui doit être ; car si la figure de la surface donnée dépend de la
différence entre les courbures principales, la courbure moyenne ne dépend,
au contraire, que de la somme des mêmes courbures principales.

Il est évident que la considération des \emph{surfaces des distances} et
celle de leurs divers développemens donnent à la

\page{80}
comparaison entre la courbure de la surface donnée et celle d'une certaine
sphère ; nous dirons donc que, s'il est certain dans le sens du passage
d'Euler rapporté N° 1, qu'on ne sauroit « même comparer la courbure des
surfaces avec celle d'une sphère », il ne l'est pas moins, dans le sens qui
a été expliqué, que l'idée complètte de la manière dont la courbure est
distribuée autour du point donné, aussi bien que la connoissance de la
distance moyenne entre les points de la surface qui environnent le point de
tangence, et le plan tangent, résulte de la comparaison de la courbure de la
surface à celle de sa sphère de moyenne courbure.

N° 8 La surface sphérique, d'une part, et de l'autre, la surface dont les rayons
de principales courbures sont égaux, mais de signes opposés, comprennent
entr'elles toutes les figures possibles de la surface courbe.

Dans ce qui précède, on s'est uniquement occupé du cas où les deux rayons de
courbures principales sont dirigés du même côté ; et, pour plus de
simplicité, on a supposé que, quel que soit le changement des figures, la
somme des courbures principales, représentée ici par celle des \emph{lignes
de principales distances}, est toujours la même.

Après avoir également supposé, à l'égard des surfaces dont les rayons de
principales courbures sont dirigés dans des sens opposés, que, abstraction
faite du signe, la somme des \emph{lignes de principales distances} demeure
toujours la même, nous examinerons ce que deviennent alors, et la
\emph{surface des}

\note{la vue 80 s'arrête au milieu de la phrase ; le texte se poursuit à la
vue 81, hors de ce lot.}

\end{document}
